\documentclass[12pt]{article}

\usepackage[a4paper, margin=2.5cm]{geometry}
\usepackage{amsmath}
\usepackage{amssymb}
\usepackage{amsthm}
\usepackage[colorlinks=true, linkcolor=blue, citecolor=blue, urlcolor=blue]{hyperref}
\usepackage{booktabs}
\usepackage{natbib}
\usepackage{enumitem}

\newtheorem{theorem}{Theorem}
\newtheorem{proposition}[theorem]{Proposition}
\newtheorem{corollary}[theorem]{Corollary}
\newtheorem{definition}[theorem]{Definition}
\newtheorem{remark}[theorem]{Remark}
\newtheorem{lemma}[theorem]{Lemma}
\newtheorem{observation}[theorem]{Observation}

\newcommand{\Fcal}{\mathcal{F}}
\newcommand{\Fint}{\Fcal_{\mathrm{int}}}
\newcommand{\Ftot}{\Fcal_{\mathrm{tot}}}
\newcommand{\HW}{H_{\mathrm{W}}}
\newcommand{\Sadm}{\mathcal{S}_{\mathrm{adm}}}
\newcommand{\Mcl}{\mathcal{M}_{\mathrm{cl}}}
\newcommand{\phig}{\varphi}
\newcommand{\lbar}{\bar{\lambda}}

\title{The Closure Functional from First Principles:\\[6pt]
\large Self-Similar Permeability, the 600-Cell, and the Derived Interaction}
\author{%
  Lee Smart\\[4pt]
  \textit{Institute of Vibrational Field Dynamics}\\[2pt]
  \texttt{contact@vibrationalfielddynamics.org}\\[2pt]
  \texttt{@vfd\_org}%
}
\date{April 2026}

\begin{document}

\maketitle

\noindent\textbf{Status:} Preprint (not peer-reviewed)

\begin{abstract}
The closure functional $\Fcal$ and the interaction functional $\Fint$ have been
the foundation of the VFD programme, but both were previously
introduced as constructions rather than derived from first principles.
This paper narrows that gap within stated structural assumptions.

Starting from the permeability axiom (the vacuum field is permeable, with a self-similar tunnelling hierarchy) --- together with
the geometric-mean closure condition, we derive $\phig$ as the unique
positive fixed point of the resulting recursion. Under the additional
structural assumptions of compact $S^3$ configuration space,
finite-subgroup closure, and maximal vertex-transitivity, the
$\phig$-determined $\pi/5$ angular scale uniquely selects the binary
icosahedral orbit within the ADE classification --- i.e., the 600-cell.
Within the class of continuous regularised soft-min functionals on
finite vertex sets satisfying sharp-limit closure, additive
composition under independent products, and the Faddeev grouping
(recursivity) axiom, the log-sum-exp functional is the Shannon/Faddeev-Khinchin choice. Product additivity alone is not sufficient (R\'enyi entropies also satisfy it), so the grouping axiom is indispensable for uniqueness; the paper states this explicitly as axiom A4. The interaction functional is
then derived from $\Fcal$ via vertex-exclusion in the two-basin
partition function: $\Fint = -\varepsilon^2\log(1 - \eta)$, where
$\eta$ is the basin overlap ratio.

A spectral decomposition of $\Fint$ motivates two retained mode sectors at nuclear distances (the other modes are assumed to be negligible at those distances, not derived as such from the $\Fint$ expansion): a \textit{phase-coherence channel}
on the Hopf fibres over the nucleon support graph (identified with
pion exchange, six shell-matched channels on $P_3$) and the
\textit{graph-Laplacian zero mode} (identified with photon exchange,
coupling $\alpha$). The resulting structural separation factor
$(6 + \alpha)$ yields a deuteron charge radius
$r_d = 2.12800$~fm, in agreement with experiment within $1\sigma$ against both CODATA anchors (CODATA 2018 $r_d = 2.12799(74)$~fm: $0.02\sigma$; CODATA 2022 $r_d = 2.12778(27)$~fm: $0.83\sigma$), improving the leading-order result by $\sim$100$\times$ (CODATA 2018) and $\sim$9$\times$ (CODATA 2022) in absolute error. The structural-separation coefficient $(6+\alpha)$ is obtained as a \emph{conditional identification} (see Proposition~\ref{thm:factor}), not a first-principles derivation from $\Fint$: the value $\pi\,\lbar_p$ of the kinematic coefficient $-g'(r_*^{(\pi)})/K_\pi$ is inherited from the half-period identification $r_\pi = \pi\,\lbar_p$ of~\citep{SmartRadius}, while the derivative/curvature identity $-g'(r_*^{(\pi)})/K_\pi = \pi\,\lbar_p$ itself is introduced here as a kinematic identification, not derived from the vertex-exclusion overlap integral of this paper.

Every step is classified as \textbf{derived}, \textbf{identified}
(structural correspondence), or \textbf{open}, so the reader can locate
exactly where the logical chain rests on each type of support.
\end{abstract}

%==================================================================
\section{Introduction}\label{sec:intro}
%==================================================================

The VFD closure programme has reported conditional numerical correspondences for particle masses at $0.014\%$
average error~\citep{SmartV}, structural constants
($\alpha^{-1} = 137 + \pi/87$)~\citep{SmartXXII}, gauge
structure~\citep{SmartVI}, and six hadron charge
radii~\citep{SmartRadius}. All results flow from the closure functional
$\Fcal$ on the 600-cell. Yet $\Fcal$ itself was introduced as a
definition~\citep{SmartXXII}, and the interaction functional $\Fint$ was
an explicit ansatz with unspecified weights~\citep{SmartXIII}.

This paper derives both within a clearly stated set of structural
axioms. The logical chain is:
\begin{enumerate}
    \item \textbf{Geometric-mean self-similar permeability} uniquely selects $\phig$ (Section~\ref{sec:phi}).
    \item Under compact $S^3$, finite-subgroup closure, and maximal vertex-transitivity, $\phig$ \textbf{selects the 600-cell} via the ADE classification (Section~\ref{sec:600cell}).
    \item The 600-cell vertex set + axioms \textbf{A1--A4} (including A3 extensivity and A4 Faddeev grouping) uniquely characterise $\Fcal$ within the class of regularised soft-min functionals (Section~\ref{sec:F}; A4 is indispensable --- product-additivity alone is also satisfied by R\'enyi entropies, so the grouping axiom is required to single out Shannon/log-sum-exp).
    \item $\Fcal$ + \textbf{vertex exclusion} yields $\Fint$ from the two-basin partition function (Section~\ref{sec:Fint}).
    \item \textbf{Heuristic spectral expansion} of $\Fint$ motivates retaining two mode sectors (phase-coherence and graph-Laplacian zero mode), identified with pion and photon exchange (Section~\ref{sec:modes}); the spectral expansion itself is heuristic, not a proved identity.
    \item The separation factor $(6 + \alpha)$ --- obtained as a conditional identification (Proposition~\ref{thm:factor}) rather than a first-principles derivation from $\Fint$ --- gives the deuteron to $0.02\sigma$ against CODATA 2018 and $0.83\sigma$ against CODATA 2022 (Section~\ref{sec:deuteron}).
\end{enumerate}

\subsection*{Epistemic classification}

Each result is marked at one of three levels:
\begin{itemize}
    \item \textbf{Derived} ($\blacktriangleright$): follows mathematically from stated axioms. Proved or provable from known theorems.
    \item \textbf{Identified} ($\triangleright$): a structural correspondence --- a quantity computed within the framework is identified with a physical observable. Numerically verified but not deduced from the axioms alone.
    \item \textbf{Open} ($\circ$): not resolved by this paper.
\end{itemize}

\subsection*{Narrative placement (closure projection)}

The deuteron radius result of this paper is an instance of the \textit{composite projection class} of the closure-projection framework introduced in Paper~XXXIII~\citep{SmartXXXIII}: a quadrature combination of the spatial projection class (boundary-resolvent and phase-coherence projections of $\Fcal$ onto an electromagnetic-probe boundary) and the interaction-mode projection from $\Fint$ derived in Sections~\ref{sec:Fint}--\ref{sec:modes}. The wave-substrate $\Fcal$ is derived from first principles in this paper itself; the observer-position is placed as a bounded, stable, self-referential constraint substructure in Paper~XXIX~\citep{SmartXXIX}, with the upstream consolidation in Paper~XXXVI~\citep{SmartXXXVI}. The closure-projection identification organises the proton, pion, and deuteron radii of the radius programme as instances of one measurement principle; it is interpretive, made in Paper~XXXIII, and does not alter the conditional identifications (H1--H4) of Sections~\ref{sec:modes}--\ref{sec:deuteron}, all of which retain the hypothesis flags imposed there.

%==================================================================
\section{Self-Similar Permeability and the Golden Ratio}\label{sec:phi}
%==================================================================

\subsection{The permeability axiom}

\begin{definition}[Permeable vacuum]\label{def:permeable}
The vacuum field admits discrete tunnelling channels between
constraint-satisfying configurations. Let $d_1^2$ be the squared
distance to the nearest-neighbour (NN) configuration and $d_2^2$ be the
squared distance to the next-nearest-neighbour (NNN) configuration. The
\textbf{permeability gap} is $\Delta = d_2^2 - d_1^2$.
\end{definition}

The WKB tunnelling amplitude between two configurations separated by
squared distance $d^2$ is $t \sim \exp(-C\,d^2/\sigma^2)$, where
$\sigma^2 = \hbar/m$ is the diffusion parameter. The gap $\Delta$
controls the ratio of NN to NNN tunnelling:
\begin{equation}\label{eq:tunnel_ratio}
    t_2/t_1 = \exp(-C\,\Delta/\sigma^2).
\end{equation}

\subsection{The geometric-mean condition}

\begin{proposition}[$\blacktriangleright$ Geometric-mean self-similarity selects $\phig$]\label{prop:phi}
Define $r = \Delta/d_1^2$ (the ratio of the permeability gap to the
NN squared distance). Among two-step tunnelling hierarchies
parametrised by $(d_1^2, \Delta, d_2^2)$ with $d_2^2 = d_1^2 +
\Delta$, the \textbf{geometric-mean self-similarity axiom}
\begin{equation}\label{eq:geomean}
    \Delta^2 \;=\; d_1^2 \cdot d_2^2 \qquad\Longleftrightarrow\qquad
    r \;=\; 1 \,+\, \tfrac{1}{r}
\end{equation}
has the unique positive solution $r = \phig = (1+\sqrt{5})/2$.
\end{proposition}

\begin{remark}[Scope of Proposition~\ref{prop:phi}]
The proposition does not claim that ``self-similar permeability''
selects $\phig$ in full generality; it selects $\phig$ under the
specific geometric-mean closure condition $\Delta^2 = d_1^2 d_2^2$.
Other self-similarity conditions (e.g., arithmetic mean, harmonic
mean) would select other ratios. The geometric-mean condition is
singled out as the two-step invariant compatible with a
scale-free log-tunnelling hierarchy (Remark~\ref{rem:hurwitz}), and
its selection of $\phig$ is then reinforced \textit{a posteriori}
by the maximal Diophantine non-approximability of $\phig$.
\end{remark}

\begin{proof}
The geometric-mean condition $\Delta^2 = d_1^2\,d_2^2$ with
$d_2^2 = d_1^2 + \Delta$ gives $\Delta^2 = d_1^2(\Delta + d_1^2)$.
Dividing by $d_1^4$: $r^2 = r + 1$, with unique positive root
$r = \phig$. Under the \emph{shell-normalisation hypothesis} $d_2 = 1$ (i.e., the NNN squared chordal distance is normalised to unity; this is a choice of scale used throughout the paper and not forced by the geometric-mean condition), $d_1^2 = 1/\phig^2$ and $\Delta = 1/\phig$. Different scale choices rescale $d_1, d_2, \Delta$ simultaneously without changing the ratios.
\end{proof}

\begin{remark}[Physical meaning of the geometric-mean condition]
The condition states: \textit{the ratio of the permeability gap to the NN
scale equals the ratio of the NNN scale to the gap}. Equivalently,
$d_1^2 : \Delta : d_2^2 = 1 : \phig : \phig^2$. In the tunnelling-rate
language, the log-suppression from NN to the gap and from the gap to NNN
are in the same ratio as the rates themselves --- the hierarchy is
scale-invariant.
\end{remark}

\begin{remark}[Maximal non-resonance (motivational)]\label{rem:hurwitz}
Hurwitz's theorem~\citep{Hurwitz1891} states that for any irrational $\xi$, there are \emph{infinitely many} rationals $p/q$ satisfying $|\xi - p/q| < 1/(\sqrt{5}\,q^2)$, and that the constant $\sqrt{5}$ is sharp exactly for $\xi$ lying in the $\mathrm{SL}_2(\mathbb Z)$-equivalence class of $\phig$ --- i.e., $\phig$ is ``most poorly approximable'' in the sense that it saturates the Hurwitz bound; no larger constant can replace $\sqrt{5}$ for this equivalence class. Concretely, if the tunnelling ratio $r$ were rational,
$r = p/q$, the tunnelling amplitudes at shells $n$ and $n{+}q$ would
satisfy $t_n/t_{n+q} = \exp(-C\,q\,\log r\,/\sigma^2)$ with rationally
commensurate phases, producing zero-gap resonant channels in the
linearised transfer operator. The $\phig$-ratio minimises the density of
such near-resonances: by Hurwitz's theorem, for any irrational $\xi$ there are infinitely many rationals $m/n$ with $|\xi - m/n| < 1/(\sqrt{5}\,n^2)$, and the constant $\sqrt{5}$ is sharp exactly for $\xi$ equivalent to $\phig$ under $\mathrm{SL}_2(\mathbb Z)$. The resonance-avoidance claim made here is \emph{motivational}, drawing on the sharpness of Hurwitz; converting it into a precise transfer-operator spectral-gap theorem is an identified open item.
\end{remark}

\subsection{The angular structure}

Throughout this paper, distances on the unit 3-sphere $S^3 \subset \mathbb{R}^4$ are taken to be \emph{chordal} (ambient Euclidean) distances; the chordal squared distance in terms of the geodesic angle $\theta$ is $d_{\mathrm{ch}}^2 = 2(1 - \cos\theta)$. Chordal and geodesic distances agree to $O(\theta^4)$ for small $\theta$ (the chordal squared distance is $\theta^2 - \theta^4/12 + O(\theta^6)$, the geodesic squared distance is $\theta^2$), and the sharp-limit axiom A2 of Section~\ref{sec:F} uses the same chordal cost as the soft-min construction~\eqref{eq:softmin_class} (the two conventions are therefore consistent with each other; see A2 and Theorem~\ref{thm:F}).
From $d_1^2 = 1/\phig^2$ (under the shell-normalisation $d_2 = 1$ of Proposition~\ref{prop:phi}):
\begin{equation}\label{eq:angle}
    \cos\theta_1 = 1 - \frac{1}{2\phig^2} = \frac{\phig}{2}
    \qquad\Longrightarrow\qquad
    \theta_1 = \frac{\pi}{5} = 36^\circ.
\end{equation}

($\blacktriangleright$ Derived.) The fundamental angular separation is the
decagonal angle $\pi/5$, forced by the geometric-mean permeability axiom applied under the chordal $S^3$-embedding convention above. The
complete set of pairwise cosines consistent with this angle and
vertex-transitivity is
$\{0, \pm 1/(2\phig), \pm 1/2, \pm \phig/2, \pm 1\}$ --- all elements
of $\tfrac{1}{2}\mathbb{Z}[\phig]$.

%==================================================================
\section{The 600-Cell from the ADE Classification}\label{sec:600cell}
%==================================================================

\subsection{Finite subgroups of SU(2)}

The configuration space $S^3 \cong \mathrm{SU}(2)$ admits a classified
set of finite subgroups (the ADE series):

\begin{table}[ht]
\centering
\caption{Finite subgroups of $\mathrm{SU}(2)$ and their NN angular separations.}
\label{tab:ade}
\begin{tabular}{@{}llccl@{}}
\toprule
Group & Type & Order & NN angle & NN degree \\
\midrule
$\mathbb{Z}_n$ & Cyclic & $n$ & $2\pi/n$ & 2 \\
$2D_n$ & Binary dihedral & $4n$ & $\pi/n$ & 2 \\
$2T$ & Binary tetrahedral & 24 & $\pi/3$ & 8 \\
$2O$ & Binary octahedral & 48 & $\pi/4$ & 6 \\
$2I$ & Binary icosahedral & 120 & $\pi/5$ & 12 \\
\bottomrule
\end{tabular}
\end{table}

\begin{theorem}[$\blacktriangleright$ The $\pi/5$ angular constraint selects $2I$ within ADE]\label{thm:600cell}
Adopt the following structural assumptions:
\begin{enumerate}
    \item[(H1)] the configuration space is $S^3 \cong \mathrm{SU}(2)$;
    \item[(H2)] the closure symmetry is a finite subgroup $G \leq \mathrm{SU}(2)$ (ensuring a discrete, vertex-transitive orbit);
    \item[(H3)] the nearest-neighbour angular separation is $\pi/5$ (as selected by Proposition~\ref{prop:phi} under geometric-mean self-similarity);
    \item[(H4)] $G$ is maximal (not a proper subgroup of another finite subgroup of $\mathrm{SU}(2)$ satisfying (H3)).
\end{enumerate}
Then, within the ADE classification of finite subgroups of
$\mathrm{SU}(2)$, $G$ is uniquely the binary icosahedral group $2I$,
whose 120 elements on $S^3$ are the vertices of the 600-cell.
\end{theorem}

\begin{remark}[Status of the structural assumptions]
Assumptions (H1) and (H2) are genuine modelling inputs, not consequences
of the permeability axiom. (H1) places the theory on a compact,
simply-connected 3-manifold --- connected to the open question of why
spacetime has $3+1$ dimensions (see Section~\ref{sec:chain}). (H2)
restricts to discrete symmetry groups, appropriate for a tunnelling
hierarchy with sharply localised configurations. (H4) selects the
largest group compatible with (H3), motivated \textit{a posteriori} by
maximal non-resonance (Remark~\ref{rem:hurwitz}). Within this stated
class of models, the theorem is an unconditional classification
statement.
\end{remark}

\begin{proof}
The finite subgroups of $\mathrm{SU}(2)$ are classified (ADE
theorem~\citep{McKay1980,Klein1884}): cyclic $\mathbb{Z}_n$, binary
dihedral $2D_n$, binary tetrahedral $2T$, binary octahedral $2O$, and
binary icosahedral $2I$. No others exist. The NN angle formula for the
binary polyhedral groups is $\theta_{\mathrm{NN}} = \pi/p$, where $p$ is
the vertex valence of the associated Platonic solid $\{3,p\}$. For the
icosahedron $\{3,5\}$: $p = 5$, giving $\theta_{\mathrm{NN}} = \pi/5$.
For the binary dihedral groups, $\theta_{\mathrm{NN}} = \pi/n$, so
$2D_5$ also has $\pi/5$; but $|2D_5| = 20 < 120 = |2I|$, and
$2D_5 < 2I$ is a proper subgroup. For cyclic groups,
$\theta_{\mathrm{NN}} = 2\pi/n = \pi/5$ requires $n = 10$, but
$|\mathbb{Z}_{10}| = 10 < 120$ and $\mathbb{Z}_{10} < 2I$. The groups
$2T$ and $2O$ have NN angles $\pi/3$ and $\pi/4$ respectively, and do
not contain $\pi/5$ as an NN separation. Exhausting the ADE list, $2I$
is the unique maximal finite subgroup of $\mathrm{SU}(2)$ with
$\theta_{\mathrm{NN}} = \pi/5$.
\end{proof}

\begin{remark}[$\triangleright$ Classification note (assumed, not proved)]\label{cor:phi_determined}
Invoking the standard classification of finite $H_4$-invariant configurations on $S^3$ (cited below), the 600-cell vertex set is the \emph{unique $H_4$-invariant} 120-vertex orbit of $2I$ in $S^3$; the 600-cell is therefore maximal among finite $H_4$-invariant subsets satisfying the pairwise-inner-product constraint, in the sense that no strict $H_4$-invariant superset preserves all pairwise inner products in $\tfrac{1}{2}\mathbb{Z}[\phig]$. The distance spectrum consists of 8 shells, all $\phig$-determined (Table~\ref{tab:shells}). Note that ``$H_4$-invariant'' is \emph{not} synonymous with ``closed under $2I$-action'': $H_4$ is the Coxeter (reflection) group of the 600-cell, while $2I$ is its binary-icosahedral rotation subgroup; we use $H_4$-invariance in its reflection-group sense. This is invoked here as an imported classification fact, not re-derived in the present paper; see the discussion below.
\end{remark}

\noindent The maximality statement under the $H_4$-invariance hypothesis is asserted on the basis of the standard classification of $H_4$-invariant configurations on $S^3$ (the 120-vertex orbit of $2I$ is the unique vertex-transitive polytope on $S^3$ whose root system is $H_4$, and $H_4$ is the unique non-crystallographic Coxeter group in dimension 4~\citep[Ch.~8]{ConwaySloane1999}); a full classification proof is not supplied here. Dropping $H_4$-invariance, the stronger statement ``unique maximal finite subset of $S^3$ whose pairwise inner products all lie in $\tfrac{1}{2}\mathbb{Z}[\phig]$'' is \emph{not} established here; it would require a separate classification result without the invariance hypothesis, which we do not supply.

\begin{table}[ht]
\centering
\caption{Distance shells of the 600-cell.}
\label{tab:shells}
\begin{tabular}{@{}rllrl@{}}
\toprule
Shell & $d^2$ & Angle & Degree & $\cos(\text{angle})$ \\
\midrule
1 & $2 - \phig$ & $36^\circ$ & 12 & $\phig/2$ \\
2 & $1$ & $60^\circ$ & 20 & $1/2$ \\
3 & $3 - \phig$ & $72^\circ$ & 12 & $1/(2\phig)$ \\
4 & $2$ & $90^\circ$ & 30 & $0$ \\
5 & $1 + \phig$ & $108^\circ$ & 12 & $-1/(2\phig)$ \\
6 & $3$ & $120^\circ$ & 20 & $-1/2$ \\
7 & $2 + \phig$ & $144^\circ$ & 12 & $-\phig/2$ \\
8 & $4$ & $180^\circ$ & 1 & $-1$ \\
\bottomrule
\end{tabular}
\end{table}

\begin{remark}[$\phig^2$-chains]
The self-similar permeability condition (Proposition~\ref{prop:phi})
holds exactly for the first two shells: $d_1^2 : \Delta : d_2^2 =
1/\phig^2 : 1/\phig : 1$. The full 8-shell spectrum organises into
$\phig^2$-chains: shells $1 \to 2 \to 5$ form a geometric progression
with ratio $\phig^2$ (squared distances $2{-}\phig, 1, 1{+}\phig$), as
do shells $3 \to 7$ ($3{-}\phig, 2{+}\phig$). The integer shells (2, 4,
6, 8) form an arithmetic progression. The self-similarity is exact at
the fundamental level and organises the higher shells into
$\phig$-determined families.
\end{remark}

%==================================================================
\section{The Closure Functional: Uniqueness}\label{sec:F}
%==================================================================

Given the 600-cell vertex set $V = \{v_1, \ldots, v_{120}\} \subset S^3$,
we derive the closure functional from four axioms.

\subsection{Axioms}

\begin{enumerate}
    \item[\textbf{A1.}] \textbf{Continuity.} $\Fcal_\varepsilon^V : M \to \mathbb{R}$ is continuous for each $\varepsilon > 0$ and each compact Riemannian manifold $M$ equipped with a finite vertex set $V \subset M$.
    \item[\textbf{A2.}] \textbf{Sharp limit.} $\Fcal_\varepsilon^V(x) \to \tfrac{1}{2}\min_{v \in V} c(x, v)$ as $\varepsilon \to 0$, where $c(x, v) = |x - v|^2$ is the squared chordal distance when $M$ is embedded in a Euclidean ambient space (so that on the unit $S^3 \subset \mathbb{R}^4$ used in this paper, $c(x, v) = 2(1 - \cos\theta_{xv})$ agrees with the squared geodesic distance $\theta_{xv}^2$ to $O(\theta^4)$). A2 thus fixes the cost function used to select the nearest vertex; geodesic and chordal squared distances select the same nearest vertex on $S^3$, so the nearest-vertex selection claim of A2 is unambiguous, but the explicit cost used in the soft-min class~\eqref{eq:softmin_class} is the chordal form throughout.
    \item[\textbf{A3.}] \textbf{Extensivity under independent composition.} For independent closure systems $(M_1, V_1)$ and $(M_2, V_2)$, the product system $(M_1 \times M_2, V_1 \times V_2)$ satisfies
    \begin{equation}\label{eq:A3}
        \Fcal_\varepsilon^{V_1 \times V_2}(x_1, x_2) = \Fcal_\varepsilon^{V_1}(x_1) + \Fcal_\varepsilon^{V_2}(x_2)
        \qquad \forall\,(x_1, x_2) \in M_1 \times M_2.
    \end{equation}
    \item[\textbf{A4.}] \textbf{Grouping (Faddeev recursivity).} The regulariser $R$ defining the soft-min class~\eqref{eq:softmin_class} satisfies the Faddeev grouping relation: for any probability vector $p = (p_1, \ldots, p_n) \in \Delta_n$ with $n \geq 3$ and any partition of the first two components into a compound,
    \begin{equation}\label{eq:A4}
        R(p_1, p_2, p_3, \ldots, p_n) \;=\; R(p_1 + p_2, p_3, \ldots, p_n) \;+\; (p_1 + p_2)\,R\!\left(\frac{p_1}{p_1 + p_2},\,\frac{p_2}{p_1 + p_2}\right),
    \end{equation}
    together with the symmetry and continuity conditions of Faddeev~\citep{Faddeev1956}.
\end{enumerate}

Axiom~A2 says the functional selects the nearest vertex in the sharp
limit --- it is a closure functional. Axiom~A3 is formulated at the
level of the \textit{family} $\{\Fcal_\varepsilon^V\}_V$ of functionals
indexed by the vertex set: two independent closure systems, living on
distinct configuration manifolds with disjoint vertex sets, must
compose additively on their product manifold. This is the analogue of
extensivity in thermodynamics: independent subsystems contribute
independently to the total free energy. Axiom~A4 is the Faddeev
grouping (recursivity) axiom in its standard form; it is required in
addition to product additivity because product additivity alone is
satisfied by the R\'enyi entropies $R_\alpha(p) = (1-\alpha)^{-1}\log\sum_i p_i^\alpha$ for every $\alpha > 0$, $\alpha \neq 1$ (all of which satisfy $R_\alpha(p \otimes q) = R_\alpha(p) + R_\alpha(q)$; the limit $\alpha \to 1$ recovers Shannon entropy), so A3 alone does not uniquely single out Shannon entropy; the grouping axiom A4 is the classical additional ingredient for uniqueness~\citep{Faddeev1956}. The single-manifold functional $\Fcal_\varepsilon^V$ on the 600-cell is then the restriction of the family to the fixed choice $(M, V) = (S^3, 2I)$.

\subsection{The variational representation}

\begin{theorem}[$\blacktriangleright$ Uniqueness of the closure functional under A1--A4]\label{thm:F}
Consider the class $\mathcal{C}$ of continuous regularised soft-min
functionals
\begin{equation}\label{eq:softmin_class}
    \Fcal_\varepsilon^V(x) \;=\; \mathrm{opt}_{p \in \Delta_{|V|}}
    \Bigl[\sum_i p_i\,c_i(x) \,+\, \varepsilon^2 R(p)\Bigr],
\end{equation}
with cost $c_i(x) = \tfrac{1}{2}|x - v_i|^2$ (half the squared chordal distance; the factor $\tfrac{1}{2}$ is a normalisation convention carried through the Gaussian kernels that appear later. Here $|x - v_i|^2$ is the ambient Euclidean squared distance in the embedding of the unit $S^3 \subset \mathbb{R}^4$, equal to $2(1-\cos\theta)$ in terms of the geodesic angle, and agrees with the squared geodesic distance to $O(\theta^4)$), regulariser $R :
\Delta_{|V|} \to \mathbb{R}$ continuous, permutation-symmetric, and
normalised so that $R(\delta_i) = 0$ at every vertex of the simplex,
and $\mathrm{opt} \in \{\min, \max\}$. Within $\mathcal{C}$, the
unique functional satisfying axioms A1--A4 (up to reparametrisation of
$\varepsilon$ and additive constants) is:
\begin{equation}\label{eq:F}
    \Fcal(x) \;=\; -\varepsilon^2\log\!\left(\sum_{i=1}^{120}
    \exp\!\left(-\frac{|x - v_i|^2}{2\varepsilon^2}\right)\right)
\end{equation}
Equivalently:
\begin{equation}\label{eq:F_var}
    \Fcal(x) \;=\; \min_{p\,\in\,\Delta_{120}}\left[
    \sum_i p_i\,\frac{|x - v_i|^2}{2} + \varepsilon^2\sum_i p_i\log p_i
    \right]
\end{equation}
where $\Delta_{120}$ is the probability simplex and the minimiser is the
Gibbs distribution $p_i^*(x) = \exp(-|x - v_i|^2/2\varepsilon^2) /
Z(x)$.
\end{theorem}

\begin{proof}
Write $\Fcal_\varepsilon = \operatorname{opt}_p[\sum_i p_i\,c_i(x)
+ \varepsilon^2 R(p)]$ for some regulariser $R$ and cost
$c_i(x) = |x - v_i|^2/2$ (squared chordal distance on the unit $S^3$). Extensivity (A3) requires $R(p \otimes q) = R(p) + R(q)$ for product distributions. Grouping (A4) together with continuity and permutation symmetry is the Faddeev axiom system; by the Faddeev--Khinchin uniqueness theorem~\citep{Faddeev1956}, the unique continuous permutation-symmetric functional on probability vectors satisfying both product additivity (A3) and grouping (A4) with $R(\delta_i) = 0$ is Shannon entropy, $R(p) = \beta\sum_i p_i\log p_i$ for some $\beta > 0$ (the constant $\beta$ is the Shannon base-of-logarithm parameter). Product additivity alone is not sufficient --- it is also satisfied by every R\'enyi entropy --- so A4 is indispensable for selecting Shannon entropy uniquely (see the discussion after A4 above). The sharp limit (A2) forces $\operatorname{opt} = \min$
(the entropy penalty vanishes as $\varepsilon \to 0$, collapsing $p^*$
onto the cost minimiser). The Shannon-base parameter $\beta$ is absorbed into the
definition of $\varepsilon$. Substituting the Gibbs minimiser into the
variational form recovers~\eqref{eq:F}. This is the Gibbs variational
principle~\citep{ShoreJohnson1980}.
\end{proof}

\begin{remark}[Why not other smooth approximations?]
The $p$-norm approximation $(\sum |x - v_i|^{-p})^{-1/p}$ fails
smoothness at vertices. The Moreau envelope $\min_y[f(y) +
|x-y|^2/2\varepsilon]$ achieves only $C^{1,1}$, not $C^\infty$. Polynomial
smoothings fail extensivity. The log-sum-exp is the \textit{only}
functional in class $\mathcal{C}$ satisfying A1--A4.
\end{remark}

\subsection{Properties}

The following are consequences, not axioms:

\begin{remark}[$\blacktriangleright / \triangleright$ Expected properties of $\Fcal$ (proof sketches)]\label{prop:properties}
The closure functional~\eqref{eq:F} has the following properties (those marked $\blacktriangleright$ admit direct proof from the log-sum-exp form, those marked $\triangleright$ are stated here as expected properties with proof-sketch level support): (1)~$C^\infty$ smoothness for $\varepsilon > 0$~$\blacktriangleright$ (from differentiation of the log-sum-exp form); (2)~$H_4$ invariance (symmetry of the vertex set under the reflection group of the 600-cell)~$\blacktriangleright$ (from $H_4$-invariance of the sum in~\eqref{eq:F}); (3)~vertex minimality ($\Fcal(v_i) \leq \Fcal(x)$ for $\varepsilon^2\log 120 < d_{\min}^2/2$, where $d_{\min}$ is the nearest-neighbour chordal distance)~$\triangleright$ (proof-sketch: follows from the sharp-limit axiom A2 by continuity at $\varepsilon \to 0$; a full quantitative proof would need the explicit Hessian estimate); (4)~monotonicity (if $|x - v_i| \geq |y - v_i|$ for all $i$, then $\Fcal(x) \geq \Fcal(y)$)~$\blacktriangleright$ (direct: the log-sum-exp is a monotone increasing function of each squared distance when all distances are increased); (5)~uniform Hessian at vertices (by vertex-transitivity)~$\blacktriangleright$ (from $H_4$-invariance of $\Fcal$ and the vertex-transitivity of the 600-cell vertex orbit); (6)~barrier height between adjacent vertices controlled by $\varepsilon$~$\triangleright$ (proof-sketch: follows from the sharp-limit behaviour as $\varepsilon \to 0$ in which the barrier tends to $\tfrac{1}{2}d_{\mathrm{mid}}^2$ where $d_{\mathrm{mid}}$ is the distance from the midpoint to the nearest vertex; a full proof with explicit $\varepsilon$-dependence is not supplied here).
\end{remark}

%==================================================================
\section{The Interaction Functional: Vertex Exclusion}\label{sec:Fint}
%==================================================================

\subsection{The two-basin partition function}

When two states $\psi_1, \psi_2$ occupy basins on the same 600-cell
landscape, the \textit{unrestricted} joint partition function
$Z(\psi_1)\,Z(\psi_2)$ factors and gives zero interaction. Interaction
arises from the constraint that two states cannot simultaneously
occupy the same vertex --- the combinatorial analogue of the Pauli
exclusion principle.

\begin{definition}[$\blacktriangleright$ Exclusion-constrained partition function]\label{def:Z2}
\begin{equation}\label{eq:Z2}
    Z_2(\psi_1, \psi_2) = \sum_{\substack{i,j=1 \\ i \neq j}}^{120}
    \exp\!\left(-\frac{|\psi_1 - v_i|^2 + |\psi_2 - v_j|^2}{2\varepsilon^2}\right)
\end{equation}
\end{definition}

\begin{proposition}[$\blacktriangleright$ Factorisation failure]\label{prop:Z2}
$Z_2 = Z(\psi_1)\,Z(\psi_2) - D(\psi_1, \psi_2)$, where the
\textbf{diagonal sum} is:
\begin{equation}\label{eq:D}
    D(\psi_1, \psi_2) = \sum_{i=1}^{120}
    \exp\!\left(-\frac{|\psi_1 - v_i|^2 + |\psi_2 - v_i|^2}{2\varepsilon^2}\right)
    = \exp\!\left(-\frac{|\psi_1 - \psi_2|^2}{4\varepsilon^2}\right)\,
    \widetilde{Z}(\bar{\psi})
\end{equation}
where $\bar{\psi} = (\psi_1 + \psi_2)/2$ and
$\widetilde{Z}(\bar{\psi}) = \sum_i \exp(-|v_i - \bar{\psi}|^2/\varepsilon^2)$
is the midpoint partition function at halved temperature.
\end{proposition}

\begin{proof}
Complete the square:
$|\psi_1 - v_i|^2 + |\psi_2 - v_i|^2 = 2|v_i - \bar{\psi}|^2 + \tfrac{1}{2}|\psi_1 - \psi_2|^2$.
\end{proof}

\subsection{The derived interaction functional}

\begin{theorem}[$\blacktriangleright$ Interaction from exclusion]\label{thm:Fint}
The total two-state closure functional is $\Ftot = -\varepsilon^2\log Z_2$.
The interaction functional
$\Fint = \Ftot - \Fcal(\psi_1) - \Fcal(\psi_2)$ is:
\begin{equation}\label{eq:Fint}
    \Fint(\psi_1, \psi_2) = -\varepsilon^2\log(1 - \eta)
    = \varepsilon^2\sum_{n=1}^{\infty}\frac{\eta^n}{n}
\end{equation}
where $\eta = D(\psi_1,\psi_2)/[Z(\psi_1)\,Z(\psi_2)]$ is the
\textbf{overlap ratio}.
\end{theorem}

\begin{proposition}[$\blacktriangleright$ Properties of $\Fint$]\label{prop:Fint_props}
\begin{enumerate}
    \item \textbf{Repulsive}: $\Fint \geq 0$ for all $\varepsilon > 0$ (since $\eta \in [0,1)$, so $-\log(1 - \eta) \geq 0$).
    \item \textbf{Cluster decomposition}: for $|\psi_1 - \psi_2| \gg \varepsilon$,
    $\Fint \leq C\,\varepsilon^2\exp(-|\psi_1 - \psi_2|^2/4\varepsilon^2) \to 0$.
    \item \textbf{Hard core (sharp-localisation limit)}: For any fixed
    $\varepsilon > 0$, $\Fint$ is bounded above by
    $-\varepsilon^2 \log(1 - \|\mathbf{p}^*\|_2^2)$ and attains its
    maximum at $\psi_1 = \psi_2$, but remains finite. Divergence
    $\Fint \to +\infty$ occurs only in the coordinated limit
    $\varepsilon \to 0$ with $\psi_1, \psi_2 \to v_i$ at a common
    vertex (so that $\eta \to 1$).
\end{enumerate}
\end{proposition}

\begin{proof}
\textit{(1) $\eta \in [0,1)$.}\ Define the vertex-indexed vectors
$a_i = \exp(-|\psi_1 - v_i|^2/2\varepsilon^2)$ and
$b_i = \exp(-|\psi_2 - v_i|^2/2\varepsilon^2)$. Then
$Z(\psi_1) = \sum_i a_i$, $Z(\psi_2) = \sum_i b_i$, and
$D(\psi_1,\psi_2) = \sum_i a_i b_i$. By Cauchy--Schwarz in
$\ell^2(\{1,\ldots,120\})$,
\begin{equation}\label{eq:CS}
    D(\psi_1,\psi_2) = \sum_i a_i b_i
    \leq \Bigl(\sum_i a_i^2\Bigr)^{1/2}\Bigl(\sum_i b_i^2\Bigr)^{1/2}
    \leq \Bigl(\sum_i a_i\Bigr)\Bigl(\sum_i b_i\Bigr) = Z(\psi_1)Z(\psi_2),
\end{equation}
using $\sum_i x_i^2 \leq (\sum_i x_i)^2$ for nonnegative $x_i$. Hence
$\eta \leq 1$. Equality holds iff $a_i \propto b_i$ \textit{and} the
support reduces to a single vertex, i.e.\ $\psi_1 = \psi_2$ \textit{and}
$\varepsilon \to 0$. For distinct basins or $\varepsilon > 0$,
$\eta < 1$ strictly, so $-\log(1 - \eta) > 0$ is well-defined and
nonnegative.

\textit{(2) Cluster decomposition.}\ From Proposition~\ref{prop:Z2},
$D(\psi_1,\psi_2) = \exp(-|\psi_1 - \psi_2|^2/4\varepsilon^2)
\widetilde{Z}(\bar\psi)$ with $\widetilde{Z}(\bar\psi) \leq 120$. Hence
$\eta \leq 120\,\exp(-|\psi_1 - \psi_2|^2/4\varepsilon^2) /
[Z(\psi_1)Z(\psi_2)]$. For states near basin centres $Z(\psi_k) \geq 1$,
giving $\Fint \leq C\varepsilon^2 \exp(-|\psi_1 - \psi_2|^2/4\varepsilon^2)$
with $C = 120$.

\textit{(3) Hard core.}\ As $\psi_1 \to \psi_2$, the exponential factor
$\exp(-|\psi_1 - \psi_2|^2/4\varepsilon^2) \to 1$, and
$D/(Z_1 Z_2) \to \sum_i (p_i^*(\psi_1))^2 \equiv \|\mathbf{p}^*\|_2^2$,
which is strictly positive and equals 1 iff $\mathbf{p}^*$ is a
delta-distribution on a single vertex (i.e., $\varepsilon \to 0$). For
finite $\varepsilon$, $\|\mathbf{p}^*\|_2^2 < 1$. When also $\psi_1$
approaches a specific vertex $v_i$ with $\varepsilon \to 0$ in a
coordinated limit, $\|\mathbf{p}^*\|_2^2 \to 1$ and $\Fint \to +\infty$.
\end{proof}

\begin{remark}[Relationship to Paper~XIII]
Paper~XIII's ansatz~\citep{SmartXIII} postulated four interaction
components with unspecified weights. The derived $\Fint = -\varepsilon^2\log(1-\eta)$ of this paper replaces those four components with a single repulsive overlap-based functional on the base 600-cell vertex graph, with no free parameters beyond the regulariser $\varepsilon$. The proximity, alignment, and compatibility terms of Paper~XIII's ansatz are naturally carried by the geometric information in the overlap ratio $\eta$; the binding term, however, is \emph{not} reproduced by the repulsive $\Fint$ on its own --- binding in the present framework arises from the separate phase-bundle lift (Proposition~\ref{thm:factor}'s pion-sector binding-minimum hypothesis (H4)), not from vertex exclusion. A full replacement of Paper~XIII's binding content by a derivation from a joint base+fibre construction is an identified open item.
\end{remark}

%==================================================================
\section{Mode Decomposition of the Interaction}\label{sec:modes}
%==================================================================

\subsection{Spectral expansion}

Let $\{\lambda_k, \phi_k\}$ be the eigenvalue--eigenfunction pairs of
the 600-cell graph Laplacian $\mathbf{L}$. The \textbf{Gaussian
spectral weight} of mode $k$ at vertex $v_a$ is:
\begin{equation}\label{eq:wk}
    w_k(v_a;\varepsilon) = \frac{1}{Z(v_a)}\sum_{i=1}^{120}
    \exp\!\left(-\frac{|v_a - v_i|^2}{2\varepsilon^2}\right)\phi_k(v_i)
\end{equation}

The 120 eigenfunctions $\{\phi_k\}_{k=0}^{119}$ of $\mathbf{L}$ form a
complete orthonormal basis of $L^2(V) \cong \mathbb{R}^{120}$:
$\sum_{k=0}^{119} \phi_k(v_i)\phi_k(v_j) = \delta_{ij}$. For states
localised at vertices $v_a, v_b$, we expand the Gaussian kernel
$K_\varepsilon(v,w) = \exp(-|v-w|^2/2\varepsilon^2)$ in this basis on
$V$, yielding the interaction to leading order in $\eta$:
\begin{equation}\label{eq:Fint_spectral}
    \Fint \approx \varepsilon^2\,\eta =
    \varepsilon^2\exp\!\left(-\frac{|v_a - v_b|^2}{4\varepsilon^2}\right)
    \sum_{k=0}^{119} w_k(v_a;\varepsilon)\,w_k(v_b;\varepsilon)
\end{equation}
where the parallelogram identity $|v_a - v_i|^2 + |v_b - v_i|^2 = \tfrac{1}{2}|v_a - v_b|^2 + 2\,|v_i - \tfrac{1}{2}(v_a+v_b)|^2$ has been used to factor the bilinear exponential $\exp(-|v_a - v_b|^2/4\varepsilon^2)$ out of the overlap $\eta$, and the residual sum $\sum_i \exp(-|v_i - m|^2/\varepsilon^2)$ (with $m = \tfrac{1}{2}(v_a+v_b)$) has been re-expressed heuristically in the Laplacian eigenbasis $\{\phi_k\}$ with weights $w_k(v_a;\varepsilon),\,w_k(v_b;\varepsilon)$ encoding the Gaussian source profiles at $v_a, v_b$. We do not supply an exact identity between $\sum_i \exp(-|v_i - m|^2/\varepsilon^2)$ and $\sum_k w_k(v_a;\varepsilon)\,w_k(v_b;\varepsilon)$: a precise completeness-relation derivation must separate the heat-kernel prefactor from the eigenbasis weights, which we leave as an identified open item. Equation~\eqref{eq:Fint_spectral} is therefore a \emph{heuristic spectral decomposition} used to motivate the two-sector mode structure below, not a proved identity ($\circ$).

The Gaussian weights $w_k$ act as a low-pass filter: modes with large
$\lambda_k$ are exponentially suppressed (the support of
$K_\varepsilon$ is concentrated on shell-1 neighbours of $v_a$, and
oscillatory eigenfunctions average to zero against a localised kernel),
and the dominant contributions come from the lowest-lying modes.

\subsection{Two surviving mode classes at nuclear distances}\label{sec:two_modes}

At inter-nucleon separations ($r \sim$ few fm), two classes of modes
contribute to $\Fint$. They live in \textit{different operator spaces}
and must be carefully distinguished.

\begin{definition}[Mode classification]\label{def:mode_classes}
\begin{enumerate}
    \item The \textbf{fibre phase-coherence channel} is the flat
    direction of the closure functional along the Hopf-fibre phase
    coordinate $\phi$ at each support vertex. This channel lives on
    the tangent bundle $TS^1$ over the vertex set --- the phase
    degree of freedom attached to each basin. We identify this
    channel with \textit{pion exchange} at nuclear distances
    ($\triangleright$; see Section~\ref{sec:fibre_hessian}).

    \item The \textbf{graph-Laplacian zero mode} is the zero
    eigenvalue ($\lambda = 0$, multiplicity~1) of the 600-cell graph
    Laplacian $\mathbf{L}$. This operator acts on $L^2(V)$ ---
    functions on the vertex set. The zero mode has exact closure
    ($\mathbf{L}\phi_0 = 0$) and propagates with zero mass gap
    ($\phi_0 = 1/\sqrt{120}$). We identify this mode with
    \textit{photon exchange} with coupling $\alpha$
    ($\triangleright$; see Section~\ref{sec:photon_exchange} and
    Papers~VI~\citep{SmartVI} and~XXII~\citep{SmartXXII}).
\end{enumerate}
\end{definition}

The two channels live in \textbf{different operator spaces}: the
fibre phase-coherence channel is a lift of $\Fcal$ along
phase-bundle coordinates, while the graph-Laplacian zero mode is an
eigenfunction on $L^2(V)$. Because they act on different spaces they are not orthogonal in a common Hilbert space. We \emph{assume} (not prove here) that the two contributions enter the total free energy as separate additive pieces on their respective operator spaces, with no cross-coupling at the order considered; deriving this additive decomposition from a joint operator-space construction (e.g., a fibred Hilbert-space completion) is an identified open item ($\circ$).

The remaining modes are \emph{assumed to be negligible at nuclear distances} (not derived here from the $\Fint$ expansion): the massive
gauge bosons ($m_W, m_Z \sim 80$--$91$~GeV) have Compton wavelengths
$\lbar_W \sim 0.002$~fm $\ll r_{\mathrm{str}}$, suppressing their
propagators by $\exp(-m_W r_{\mathrm{str}}) \sim e^{-1600}$. The
irrational-sector modes are exactly decoupled from the integer sector by
the representation-theoretic theorem of Paper~XXII.

\subsection{Fibre phase-coherence channels and the pion correspondence}\label{sec:fibre_hessian}

We now make the ``fibre phase-coherence channel'' of
Definition~\ref{def:mode_classes}(1) precise, and state carefully the
correspondence with pion exchange. This section provides a
\textit{structural} account of the six-channel counting;
Proposition~\ref{prop:pion} below is therefore labelled
identified ($\triangleright$) rather than derived.

\subsubsection*{Hopf structure of the configuration manifold}

The configuration manifold $S^3 \cong \mathrm{SU}(2)$ is the total space
of the Hopf fibration
\begin{equation}\label{eq:hopf}
    U(1) \hookrightarrow S^3 \xrightarrow{\,\pi_H\,} S^2,
\end{equation}
with fibre action $R_\theta : x \mapsto x \cdot e^{i\theta}$ (right
quaternion multiplication by $e^{i\theta} \in U(1) \subset \mathrm{SU}(2)$,
$\theta \in [0, 2\pi)$). For each vertex $v_i \in V$, the Hopf fibre
through $v_i$ is $F_i = \{v_i e^{i\theta} : \theta \in [0, 2\pi)\}
\subset S^3$; since $2I$ is discrete, $F_i \cap V = \{v_i\}$. Each
basin centre therefore lies on its own Hopf fibre, giving the vertex
set $V$ a natural $U(1)$-bundle structure.

\subsubsection*{Phase-lifted closure states}

The closure functional~\eqref{eq:F} assigns a free energy to a
spatial configuration $x \in S^3$. A \textit{phase-lifted closure
state} is the pair $\Psi_i = (v_i, \phi_i)$ consisting of a basin
centre $v_i \in V$ together with an internal Hopf-fibre phase
$\phi_i \in S^1 \cong F_i$. We define the lifted functional on the
extended configuration space $V \times (S^1)^{|V|}$ by
\begin{equation}\label{eq:F_gauged}
    \Fcal_{\mathrm{g}}(x; \{\phi_i\}) \;\equiv\; \Fcal(x).
\end{equation}

\begin{remark}[Status of the phase lift]\label{rem:phase_lift_status}
The lift~\eqref{eq:F_gauged} is \textit{trivial}: by construction,
$\Fcal_{\mathrm{g}}$ has no $\phi_i$-dependence because the raw closure
functional~\eqref{eq:F} depends only on Euclidean distances between
$x$ and the vertex set. Therefore any flatness of
$\Fcal_{\mathrm{g}}$ in the fibre directions is a kinematic property
of the lift, not a consequence of a nontrivial symmetry of a
physical action. In particular, Goldstone's theorem does not apply
directly to~\eqref{eq:F_gauged}, because there is no genuine
continuous symmetry of a physical action being spontaneously
broken: the phases $\phi_i$ are auxiliary coordinates with no
kinetic term. The fibre-flatness observation below is a structural
statement about the configuration space, and the identification with
pion exchange is a physical correspondence, not a derivation.
\end{remark}

\begin{observation}[Fibre-direction flatness]\label{obs:fibre_flat}
From~\eqref{eq:F_gauged},
\begin{equation}\label{eq:hess_zero}
    \frac{\partial^2 \Fcal_{\mathrm{g}}}{\partial \phi_i\,\partial \phi_j}(x; \{\phi_k\}) \;=\; 0
    \qquad \forall\, i, j \in \{1,\ldots,120\}.
\end{equation}
The Hessian of the lifted functional restricted to the fibre
directions is identically zero. The fibre bundle
$V \times (S^1)^{|V|}$ therefore carries $|V|$ flat coordinate
directions per configuration.
\end{observation}

Observation~\ref{obs:fibre_flat} records a tautology of the trivial
lift~\eqref{eq:F_gauged}, not a dynamical result. It identifies a
natural set of $|V|$ phase coordinates at which a full field theory
(equipped with a nontrivial $U(1)^{|V|}$-invariant kinetic term)
would exhibit massless fluctuations. Within such a theory,
Goldstone's theorem would supply $|V(P_3)|$ massless modes per
nucleon upon basin localisation; the present paper does not
construct that theory.

\subsubsection*{Structural correspondence with pion exchange}

\begin{proposition}[$\triangleright$ Pion-exchange channel count, conditional]\label{prop:pion}
Adopt the structural identification (Definition~\ref{def:mode_classes}(1))
of the fibre phase-coherence channel with pion exchange at nuclear
distances, and assume the following additional structural hypothesis, \emph{not derived here} from the two-basin vertex-exclusion constraint alone: the two nucleon probes localise on \emph{disjoint} three-vertex support graphs $P_3^{(1)}, P_3^{(2)}$ with $V(P_3^{(1)}) \cap V(P_3^{(2)}) = \emptyset$. (The constraint $i \neq j$ in Definition~\ref{def:Z2} excludes coincidence of a single pair but does not by itself force disjoint support; disjointness is the additional hypothesis.) Each nucleon has support graph $P_3$ on shells
$\{2,3,4\}$, with $|V(P_3)| = 3$ occupied vertices; by
Observation~\ref{obs:fibre_flat}, the lifted closure functional has
three flat phase directions per nucleon. Under the disjoint-support hypothesis, the two nucleons' phase directions are independent coordinates on the extended configuration space.
The total pion-exchange channel count is therefore
\begin{equation}\label{eq:pion_count}
    n_\pi \;=\; 2 \times |V(P_3)| \;=\; 6.
\end{equation}
\end{proposition}

\begin{remark}[What is derived vs.\ identified in Proposition~\ref{prop:pion}]
The counting $2 \times |V(P_3)| = 6$ is a \textbf{conditional} structural
count: given the support graph $P_3$ and the \emph{assumed} disjoint-support hypothesis stated in Proposition~\ref{prop:pion}, the number of
independent flat phase coordinates in the extended configuration
space is exactly~6. The disjoint-support hypothesis is not derived from the $i \neq j$ vertex-exclusion constraint of Definition~\ref{def:Z2} alone; it is an identified open item (Section~\ref{sec:chain}). What is \textbf{identified} ($\triangleright$)
is that these six phase-coherence channels correspond to pion
exchange between the two nucleons, with coherence length $\pi\,\lbar_p$
per channel (Section~\ref{sec:pion_coherence} below). The
identification is supported by: (i) the charge-radius agreement
$r_\pi = \pi\,\lbar_p$ derived directly in the radius
paper~\citep{SmartRadius} (0.26\% error); (ii) the pion's
pseudo-Goldstone status in QCD, making its association with phase
coherence on the confinement orbit physically natural; (iii) the
numerical success of the resulting deuteron radius. It is not
derived from a dynamical action on $V \times (S^1)^{|V|}$, which
remains open.
\end{remark}

\subsection{Pion coherence length}\label{sec:pion_coherence}

Under the structural identification of Proposition~\ref{prop:pion},
each of the six phase-coherence channels contributes one coherence
length to the inter-basin structural separation. Identifying the
channel with a phase excitation on the Hopf fibre $S^1$ of angular
period $2\pi$, the natural coherence length is the half-period
(the maximum phase displacement before wrap-around at $\pi$):
\begin{equation}\label{eq:r_pi_hp}
    r_\pi \;=\; \pi\,\lbar_p \qquad (\triangleright\text{ identified}),
\end{equation}
with $\lbar_p = \hbar/(m_p c)$ setting the physical scale. Summing
over the $n_\pi = 6$ channels:
\begin{equation}\label{eq:r_pion}
    r_{\mathrm{str}}^{(\pi)} \;=\; n_\pi\,r_\pi \;=\; 6\,\pi\,\lbar_p.
\end{equation}
This matches the leading-order structural result
of~\citet{SmartRadius}, where $r_\pi = \pi\,\lbar_p$ is also
identified with the pion charge radius (experimental:
$0.659 \pm 0.004$~fm, 0.26\% error). The $\pi\,\lbar_p$ kinematic
coefficient is \textit{not} derived in the present paper from a
dynamical action on the phase bundle; it is imported from the
half-period identification in~\citep{SmartRadius} and is part of
the structural correspondence with pion exchange.

\subsection{Photon exchange: the graph-Laplacian zero mode}\label{sec:photon_exchange}

\begin{proposition}[$\triangleright$ Photon exchange]\label{prop:photon}
The $\lambda = 0$ mode of the graph Laplacian $\mathbf{L}$ propagates at
all distances with zero mass gap; its eigenfunction is the constant
mode $\phi_0 = 1/\sqrt{120}$. The coupling of this mode to each closure
basin is identified with the fine-structure constant $\alpha = 1/(137 +
\pi/87)$ (Paper~XXII, reviewed below). Under the kinematic
correspondence of Section~\ref{sec:fibre_hessian} (massless modes on
the closure orbit contribute in units of $\pi\,\lbar_p$ per channel,
weighted by coupling strength), the photon exchange contributes at
first order in $\alpha$:
\begin{equation}\label{eq:r_photon}
    r_{\mathrm{str}}^{(\gamma)} \;=\; \alpha\,\pi\,\lbar_p.
\end{equation}
\end{proposition}

\subsubsection*{Review: $\alpha^{-1} = 137 + \pi/87$ from the 600-cell
eigenvalue algebra}

For reviewer convenience, we transcribe the key arithmetic correspondence from
Paper~XXII~\citep{SmartXXII}.

\paragraph{Tree level ($137$).}
The four integer mass eigenvalues of the 600-cell closure Hamiltonian
are $\{\lambda_i\} = \{9,12,14,15\}$ (Paper~V~\citep{SmartV}). The
graph invariant
\begin{equation}\label{eq:alpha_tree}
    (\lambda_1 - 1)(\lambda_2 - 1) \;-\; 1 \;+\; \sum_{i=1}^4 \lambda_i
    \;=\; 87 \;+\; 50 \;=\; 137
\end{equation}
is exact. The number $87$ in~\eqref{eq:alpha_tree} appears in three related arithmetic identities (Paper~XXII documents explicitly that the three routes are not algebraically independent and ultimately reflect the same underlying $E_8$/2I structure):
\begin{enumerate}
    \item $(9-1)(12-1) - 1 = 8 \times 11 - 1 = 87$ \qquad (mass-eigenvalue algebra),
    \item $3(14+15) = 3 \times 29 = 87$ \qquad (Paper~V eigenvalue decomposition),
    \item $(17+19+23+29) - 1 = 87$ \qquad (upper Coxeter exponents of $E_8$).
\end{enumerate}
The coincidence $\sum(\text{upper exponents}) = 88 = (9-1)(12-1)$
connects the $E_8$ Coxeter structure to the mass eigenvalue algebra;
the total exponent sum $\sum m_i = 120$ equals the vertex count.

\paragraph{One-loop correction ($\pi/87$).}
The Coxeter element of $E_8$ (order 30) acts on the upper (conjugate)
spinor sector with phase accumulation $2\pi \times 88$. The ratio
$240/120 = 2$ records the $240$-to-$120$ two-to-one icosian/Coxeter pairing between the $E_8$ root system and the 600-cell vertex set (counting/projection pairing in the Galois-involutive sense of Paper~XXII / \texttt{cascade-algebraic-substrate}, not a topological double cover of manifolds); dividing the phase by 2 gives
$\pi \times 88$ per 600-cell cycle. Per degree of freedom
($87 = 88 - 1$, subtracting the cyclic redundancy),
\begin{equation}\label{eq:alpha_oneloop}
    \text{phase per d.o.f.} \;=\; \frac{\pi \times 88}{87} \;=\; \pi \;+\; \frac{\pi}{87}.
\end{equation}
The excess over the natural half-turn $\pi$ is $\pi/87$. Combining
with~\eqref{eq:alpha_tree}:
\begin{equation}\label{eq:alpha_full}
    \alpha^{-1} \;=\; 137 \;+\; \frac{\pi}{87} \;=\; 137.036\ldots
\end{equation}
which agrees with the CODATA value $137.035999\ldots$ to $0.81$~ppm.

\paragraph{Epistemic note.}
The pion channel count ($n_\pi = 6$) is a \textbf{conditional structural count}, obtained from the fibre--support structure (Proposition~\ref{prop:pion}) under the disjoint-support hypothesis stated in that proposition. The
photon coupling ($\alpha$) is \textbf{identified}:
equations~\eqref{eq:alpha_tree}--\eqref{eq:alpha_full} compute the
number $137 + \pi/87$ from the 600-cell eigenvalue algebra and
identify it with the physical fine-structure constant. This
identification is supported by the 0.81~ppm numerical agreement and
the three related arithmetic appearances of the tree-level term $87$ (which are not algebraically independent; see the discussion above), but the chain
from the vertex-exclusion overlap~\eqref{eq:Fint_spectral} to the
exact value of the photon coupling coefficient has not been
computed from the inter-basin overlap integral alone. Closing this
link --- deriving $\alpha$ directly from the $\Fint$ spectral
expansion rather than importing it from the Coxeter phase account
--- remains open ($\circ$).

%==================================================================
\section{The Structural Separation Factor}\label{sec:combined}
%==================================================================

\begin{proposition}[$\triangleright$\ $(6 + \alpha)$ separation factor, conditional identification]\label{thm:factor}
Under the following hypotheses:
\begin{enumerate}[label=\textup{(H\arabic*)}]
    \item Proposition~\ref{prop:pion} (six phase-coherence channels on the $P_3$ support graph under the disjoint-support hypothesis of that proposition, $\triangleright$ identified with pion exchange);
    \item Proposition~\ref{prop:photon} (graph-Laplacian zero mode $\triangleright$ identified with photon exchange at coupling $\alpha$);
    \item the kinematic identification~\eqref{eq:kinematic_id} (each massless channel contributes one half-period $\pi\,\lbar_p$ of the closure orbit, imported from~\citep{SmartRadius} and \emph{not} derived from the vertex-exclusion overlap integral of the present paper);
    \item \textbf{Pion binding minimum (assumed).} The pion-sector interaction $\Fint^{(\pi)}$ has a \emph{local minimum} at $r_*^{(\pi)} = 6\,\pi\,\lbar_p$ with strictly positive curvature $K_\pi \equiv \Fint^{(\pi)\prime\prime}(r_*^{(\pi)}) > 0$. This is an additional assumption, \emph{not} derived from the base-vertex-graph $\Fint = -\varepsilon^2\log(1-\eta)$ (which is repulsive, Proposition~\ref{prop:Fint_props}); the binding minimum is a feature of the separate phase-bundle lift, and bridging the two --- exhibiting a genuine pion binding minimum from a joint base+fibre construction --- is an identified open item.
\end{enumerate}
Under (H1)--(H4), the structural separation between two nucleon closure basins takes the form
\begin{equation}\label{eq:r_str}
    r_{\mathrm{str}} \;=\; (6 + \alpha)\,\pi\,\lbar_p \;+\; O(\alpha^2\,\pi\,\lbar_p),
\end{equation}
obtained after the kinematic coefficient $-g'(r_*^{(\pi)})/K_\pi = \pi\,\lbar_p$ is identified. Deriving that coefficient directly from the $\Fint$ spectral expansion is an identified open item ($\circ$).
\end{proposition}

\begin{proof}
Let $r$ denote the inter-basin separation. The total free energy along
the separation coordinate is
\begin{equation}\label{eq:Ftot_r}
    \Ftot(r) \;=\; \Fcal(\psi_1) \,+\, \Fcal(\psi_2) \,+\, \Fint^{(\pi)}(r) \,+\, \Fint^{(\gamma)}(r),
\end{equation}
where the two single-basin terms are $r$-independent and
$\Fint^{(\pi)}(r)$ is minimised at $r_*^{(\pi)} \equiv
r_{\mathrm{str}}^{(\pi)} = 6\,\pi\,\lbar_p$
(Prop.~\ref{prop:pion}, Eq.~\eqref{eq:r_pion}). Write the pion
contribution locally as
\begin{equation}\label{eq:pion_quad}
    \Fint^{(\pi)}(r) \;=\; \Fint^{(\pi)}(r_*^{(\pi)}) \;+\; \tfrac{1}{2}\,K_\pi\,(r - r_*^{(\pi)})^2 \;+\; O((r - r_*^{(\pi)})^3),
\end{equation}
with curvature $K_\pi \equiv \Fint^{(\pi)\prime\prime}(r_*^{(\pi)}) > 0$
(the pion channel is binding, so the Hessian is positive-definite at
its minimum). The photon contribution enters linearly at order $\alpha$:
\begin{equation}\label{eq:gamma_lin}
    \Fint^{(\gamma)}(r) \;=\; \alpha\,g(r),
\end{equation}
where $g(r) = O(1)$ is the single-photon overlap kernel evaluated from
\eqref{eq:Fint_spectral} on the $\phi_0$ mode. Substituting into
\eqref{eq:Ftot_r} and setting $\partial_r \Ftot = 0$:
\begin{equation}\label{eq:foc}
    K_\pi\,(r - r_*^{(\pi)}) \;+\; \alpha\,g'(r) \;=\; O((r-r_*^{(\pi)})^2).
\end{equation}
Evaluating $g'$ at the leading-order minimum and iterating once:
\begin{equation}\label{eq:shift}
    r_{\mathrm{str}} \;=\; r_*^{(\pi)} \;-\; \alpha\,\frac{g'(r_*^{(\pi)})}{K_\pi} \;+\; O(\alpha^2).
\end{equation}
At this point we invoke the \textbf{kinematic identification} (marked
$\triangleright$) that bridges the structural calculation to the
numerical correction:

\medskip
\noindent\fbox{\parbox{0.95\linewidth}{%
\textbf{Identification ($\triangleright$).}\ Both the pion and the
photon are massless modes on the same closure orbit of circumference
$2\pi\,\lbar_p$, integrated against the same Gaussian kernel
$K_\varepsilon$ in~\eqref{eq:Fint_spectral}. Under the kinematic
correspondence of Section~\ref{sec:fibre_hessian} (each massless
channel contributes one half-period $\pi\,\lbar_p$ of the closure
orbit, weighted by coupling strength), the first-order shift
coefficient in~\eqref{eq:shift} is
\begin{equation}\label{eq:kinematic_id}
    -\,\frac{g'(r_*^{(\pi)})}{K_\pi} \;=\; \pi\,\lbar_p.
\end{equation}
This identification is \textit{not} derived from the
vertex-exclusion overlap integral~\eqref{eq:Fint_spectral} in the
present paper; its justification is the $r_\pi = \pi\,\lbar_p$
identification inherited from~\citep{SmartRadius}. Deriving
$\pi\,\lbar_p$ directly from the $\Fint$ spectral expansion remains
open ($\circ$, Section~\ref{sec:chain}).
}}
\medskip

Substituting~\eqref{eq:kinematic_id} into~\eqref{eq:shift}:
\begin{equation}\label{eq:6plusalpha}
    r_{\mathrm{str}} \;=\; 6\,\pi\,\lbar_p \,+\, \alpha\,\pi\,\lbar_p \,+\, O(\alpha^2)
    \;=\; (6 + \alpha)\,\pi\,\lbar_p \,+\, O(\alpha^2).
\end{equation}
The $O(\alpha^2)$ neglected terms arise from the quadratic correction
in~\eqref{eq:foc}, which scales as $(\alpha/K_\pi)^2 K_\pi = O(\alpha^2)$ in the coefficient of $(r - r_*^{(\pi)})^2$. With the dimensionless expansion parameter $\alpha/n_\pi \approx 1.2 \times 10^{-3}$, these terms are \emph{expected to be small} but are \emph{not rigorously bounded here} without an explicit estimate of the hidden constant (the constant depends on $K_\pi$ and on the third derivative $g''(r_*^{(\pi)})$, neither of which is computed in the present paper); whether the $O(\alpha^2)$ correction sits below experimental precision is therefore an open item, not established.
\end{proof}

\begin{remark}[Additivity is perturbative, not exact]
The linear addition $r_{\mathrm{str}} = r_{\mathrm{str}}^{(\pi)} +
r_{\mathrm{str}}^{(\gamma)}$ holds to first order in $\alpha$ because
$\Fint^{(\gamma)}$ is a small perturbation. The $O(\alpha^2)$ terms,
while negligible ($\sim 5 \times 10^{-5}\,\pi\,\lbar_p$), are not
computed here. Whether they vanish exactly (by spectral completeness of
the integer sector) or are merely small is an open question ($\circ$).
\end{remark}

%==================================================================
\section{Application: Deuteron Charge Radius}\label{sec:deuteron}
%==================================================================

The deuteron charge radius decomposes as~\citep{SmartRadius}:
\begin{equation}\label{eq:deuteron}
    r_d^2 = r_p^2 + \langle r_n^2\rangle + \frac{r_{\mathrm{str}}^2}{4}
\end{equation}

With the conditional separation factor (Proposition~\ref{thm:factor}):
\begin{align}
    r_p &= 4\,\lbar_p = 0.84124~\text{fm}
    && (\blacktriangleright\text{ Appendix~\ref{app:proton}}) \label{eq:rp} \\
    \langle r_n^2\rangle &= -\tfrac{8}{3}\,\lbar_n^2 = -0.11762~\text{fm}^2
    && (\blacktriangleright\text{ \citep{SmartRadius}}) \label{eq:rn2} \\
    r_{\mathrm{str}} &= (6 + \alpha)\,\pi\,\lbar_p = 3.96905~\text{fm}
    && (6\text{: }\blacktriangleright;\;\alpha\text{: }\triangleright)
    \label{eq:rstr}
\end{align}
where the reduced Compton wavelengths are taken to full available precision from CODATA~\citep{CODATA2022}: $\lbar_p = \hbar/(m_p c) = 0.2103089103(24)$~fm, $\lbar_n = \hbar/(m_n c) = 0.2100194143(24)$~fm; the fine-structure constant used here is the VFD-identified value $\alpha^{-1} = 137 + \pi/87 = 137.0361103...$ (VFD value; differing from CODATA 2022's $\alpha^{-1} = 137.035999177(21)$ by $\sim 1.11 \times 10^{-4}$ in the reciprocal, i.e., $\sim 0.81$~ppm --- within the VFD programme's claimed numerical correspondence precision for $\alpha^{-1}$ (Paper~XXII), but $\sim 5300$ CODATA standard deviations apart, since CODATA's uncertainty on $\alpha^{-1}$ is at the $2.1 \times 10^{-8}$ level). The rounded values $\lbar_p = 0.21031$~fm shown above are for readability; the $0.02\sigma$ and $0.83\sigma$ headline figures use the full-precision inputs.

\begin{remark}[Supporting derivations for $r_p$ and $\langle r_n^2 \rangle$]
Equation~\eqref{eq:rp} is the boundary-resolvent result on the proton
support graph $P_3$ (Appendix~\ref{app:proton}); the factor $4$ is
\textit{structurally triply determined} --- it equals
$\mathrm{Tr}(L(P_3)) = 2|E(P_3)|/|V(P_3)| \times |V(P_3)| = 4$, equals
$\max(S_p) = 4$ (shell index), and equals the leading cumulant
coefficient $\sqrt{12 \cdot \mathrm{Tr}(L)/|V|}$ of the heat-kernel
expansion on $P_3$. Equation~\eqref{eq:rn2} follows from the dipole
mode $v_1 = [1, 0, -1]/\sqrt{2}$ of the $P_3$ Laplacian (eigenvalue
$\lambda_1 = 1$): $\langle r_n^2 \rangle = -q \,(n_{\max}^2 -
n_{\min}^2)\,\lbar_n^2$ with $q = \mathrm{Var}(\deg) = 2/9$ and
$(n_{\max}^2 - n_{\min}^2) = 16 - 4 = 12$, giving
$\langle r_n^2 \rangle = -(2/9)(12)\lbar_n^2 = -(8/3)\lbar_n^2$.
\end{remark}

\begin{proposition}[$\triangleright$ Deuteron charge radius, computed conditional on Proposition~\ref{thm:factor}]\label{prop:deuteron}
\begin{equation}\label{eq:rd}
    r_d = \sqrt{r_p^2 + \langle r_n^2\rangle +
    \frac{(6{+}\alpha)^2\pi^2\lbar_p^2}{4}}
    = 2.12800~\text{fm}
\end{equation}
\end{proposition}

\begin{proof}
$r_d^2 = 0.70768 - 0.11762 + 3.93834 = 4.52840$~fm$^2$;
$r_d = \sqrt{4.52840} = 2.12800$~fm.
\end{proof}

\begin{table}[ht]
\centering
\caption{Deuteron charge radius: leading order vs.\ conditional NLO correction (Proposition~\ref{thm:factor}), against two experimental anchors. The $\sigma$ column depends on which anchor is used; both are reported.}
\label{tab:deuteron}
\begin{tabular}{@{}lllcccc@{}}
\toprule
& Formula & Status & VFD (fm) & Expt.\ (fm) & Error & $\sigma$ \\
\midrule
LO, CODATA 2018 & $r_{\mathrm{str}} = 6\pi\lbar_p$ & $\triangleright$ &
    2.1258 & $2.12799(74)$~\citep{Pohl2016} & 0.10\% & 3.0 \\
\textbf{NLO, CODATA 2018} & $\boldsymbol{r_{\mathrm{str}} = (6{+}\alpha)\pi\lbar_p}$ &
    $\triangleright$ &
    \textbf{2.12800} & $\boldsymbol{2.12799(74)}$~\citep{Pohl2016} & \textbf{0.0006\%} &
    \textbf{0.02} \\
LO, CODATA 2022 & $r_{\mathrm{str}} = 6\pi\lbar_p$ & $\triangleright$ &
    2.1258 & $2.12778(27)$~\citep{CODATA2022} & 0.09\% & 7.4 \\
\textbf{NLO, CODATA 2022} & $\boldsymbol{r_{\mathrm{str}} = (6{+}\alpha)\pi\lbar_p}$ &
    $\triangleright$ &
    \textbf{2.12800} & $\boldsymbol{2.12778(27)}$~\citep{CODATA2022} & \textbf{0.010\%} &
    \textbf{0.83} \\
\bottomrule
\end{tabular}
\end{table}

The NLO result is within $1\sigma$ of both CODATA 2018 (based on the Pohl~\emph{et~al.} muonic deuterium measurement) and the tighter CODATA 2022 adjustment (which combines muonic and electronic data with reduced uncertainty). The $\sigma$ headline therefore depends on which anchor is adopted; the NLO formula improves on LO by a factor of $\sim$100 in absolute error against CODATA 2018 and by a factor of $\sim$9 against CODATA 2022.

%==================================================================
\section{Testable Prediction}\label{sec:prediction}
%==================================================================

\begin{quote}
\textbf{Assumption ($\triangleright$, prediction).} Assuming linear charge scaling for the photon contribution to the inter-basin separation in a nucleus with $Z$ protons,
\begin{equation}\label{eq:Zalpha}
    r_{\mathrm{str}}^{(\gamma)}(Z) = Z\,\alpha\,\pi\,\lbar_p,
\end{equation}
the per-proton correction gives $0.005$--$0.015$~fm for light nuclei. This is \emph{not} a derivation from $\Fint$; it is a charge-additive extrapolation of the $Z=1$ result of Proposition~\ref{thm:factor}, stated as a testable prediction to be checked against precision isotope-shift spectroscopy and muonic-atom measurements~\citep{Pohl2010}.
\end{quote}

\begin{table}[ht]
\centering
\caption{Predicted electromagnetic corrections.}
\label{tab:Zalpha}
\begin{tabular}{@{}lccl@{}}
\toprule
Nucleus & $Z$ & $Z\alpha$ & EM correction (fm) \\
\midrule
Deuteron & 1 & 0.00730 & 0.00482 \\
Helion & 2 & 0.01459 & 0.00964 \\
${}^6$Li & 3 & 0.02189 & 0.01446 \\
\bottomrule
\end{tabular}
\end{table}

%==================================================================
\section{Summary of the Logical Chain}\label{sec:chain}
%==================================================================

\begin{table}[ht]
\centering
\caption{Complete derivation chain with epistemic status.}
\label{tab:chain}
\begin{tabular}{@{}rlll@{}}
\toprule
Step & Result & Status & Method \\
\midrule
1 & $\phig$ from geometric-mean self-similarity & $\blacktriangleright$ & Unique positive fixed point of $r = 1 + 1/r$ \\
2 & 600-cell from $\phig$, under (H1)--(H4) & $\blacktriangleright^\ast$ & ADE classification within stated axioms \\
3 & $\Fcal$ (log-sum-exp) & $\blacktriangleright^\ast$ & Faddeev--Khinchin uniqueness within class $\mathcal{C}$ \\
4 & $\Fint$ from vertex exclusion & $\blacktriangleright$ & Partition function ratio, Cauchy--Schwarz bound \\
5a & Six flat phase-coherence channels on $P_3$ & $\blacktriangleright^\ast$ & Structural count $2 \times |V(P_3)|$ under disjoint-support hypothesis (Prop.~\ref{prop:pion}) \\
5b & Identification of 5a with pion exchange & $\triangleright$ & Correspondence with $r_\pi = \pi\,\lbar_p$~\citep{SmartRadius} \\
5c & Photon: graph-Laplacian $\lambda{=}0$ mode & $\blacktriangleright$ & Automatic from spectral graph theory \\
5d & Photon coupling $= \alpha = 1/(137 + \pi/87)$ & $\triangleright$ & $0.81$~ppm match to Paper~XXII eigenvalue algebra \\
5e & Half-period kinematic coefficient $\pi\,\lbar_p$ & $\triangleright$ & Bridge identification~\eqref{eq:kinematic_id} \\
6 & $(6{+}\alpha)$ factor & $\triangleright$ & Conditional identification under 5a--5e (Proposition~\ref{thm:factor}) \\
7 & $r_d = 2.12800$~fm & computed & $0.0006\%$ / $0.02\sigma$ vs CODATA 2018; $0.010\%$ / $0.83\sigma$ vs CODATA 2022 \\
\bottomrule
\end{tabular}
\end{table}

\noindent\footnotesize{$\blacktriangleright^\ast$ denotes ``derived
within the stated class of models or axioms'' (as distinct from
unconditional first-principles derivation).}

\subsection{What is derived ($\blacktriangleright$)}

\begin{enumerate}
    \item $\phig$ as the unique positive fixed point of the geometric-mean self-similarity recursion $r = 1 + 1/r$ (Proposition~\ref{prop:phi}).
    \item The $\pi/5$ angular scale from the permeability ratio on the round $S^3$ metric (Eq.~\ref{eq:angle}).
    \item Under (H1)--(H4), the selection of $2I$ (hence the 600-cell) within the ADE classification (Theorem~\ref{thm:600cell}): a derivation within the stated structural assumptions, not unconditional.
    \item Within the class $\mathcal{C}$ of continuous regularised soft-min functionals, the uniqueness of the log-sum-exp form $\Fcal$ satisfying A1--A4 (Theorem~\ref{thm:F}).
    \item The interaction functional $\Fint = -\varepsilon^2\log(1-\eta)$ from vertex exclusion in the two-basin partition function, with cluster decomposition, repulsion ($\eta \in [0,1)$ via Cauchy--Schwarz), and sharp-localisation hard core (Theorem~\ref{thm:Fint}, Proposition~\ref{prop:Fint_props}).
    \item The structural count of six flat phase-coherence coordinates on the deuteron configuration bundle: $n = 2 \times |V(P_3)| = 6$ (Proposition~\ref{prop:pion}, structural counting part; conditional on the disjoint-support hypothesis for the two $P_3$ probes, see the open-items list below).
    \item The graph-Laplacian zero mode as an automatic consequence of spectral graph theory ($\phi_0 = \text{constant}$).
    \item The first-order minimisation giving $r_{\mathrm{str}} = (6 + \alpha)\,\pi\,\lbar_p + O(\alpha^2)$, once the kinematic identifications are granted (Proposition~\ref{thm:factor}).
\end{enumerate}

\subsection{What is identified ($\triangleright$)}

\begin{enumerate}
    \item \textbf{Phase channels} $\leftrightarrow$ \textbf{pion exchange.} The six flat phase-coherence coordinates on the extended configuration bundle are identified with pion-exchange channels. This is a structural correspondence, not a derivation from a dynamical action (Remark~\ref{rem:phase_lift_status}); its support is the independent $r_\pi = \pi\,\lbar_p$ identification in~\citep{SmartRadius} and the numerical success of the deuteron fit.
    \item \textbf{Graph-Laplacian zero mode} $\leftrightarrow$ \textbf{photon exchange.} The $\lambda = 0$ eigenfunction is identified with photon exchange. The coupling strength $\alpha = 1/(137 + \pi/87)$ is computed in Paper~XXII from the 600-cell eigenvalue algebra (Eqs.~\ref{eq:alpha_tree}--\ref{eq:alpha_full}, 0.81~ppm agreement); the identification of this number with the physical fine-structure constant is not derived from the vertex-exclusion integral.
    \item \textbf{Half-period kinematic coefficient} $\pi\,\lbar_p$. Each massless channel contributes one half-period of the closure orbit per unit coupling (Eq.~\ref{eq:kinematic_id}); this bridges the structural count to the numerical correction and inherits its $\pi\,\lbar_p$ value from~\citep{SmartRadius}.
    \item \textbf{$Z\alpha$ scaling.} The extension to heavier nuclei assumes the photon coupling scales linearly with nuclear charge (charge-additive extrapolation, Section~\ref{sec:prediction}).
\end{enumerate}

\subsection{What is open ($\circ$)}

\begin{enumerate}
    \item \textbf{Why $S^3$?} The choice of a compact, simply-connected 3-manifold as the configuration space is assumed (H1), not derived. It is connected to the open question of why spacetime has $3+1$ dimensions.
    \item \textbf{Why a finite subgroup?} (H2) restricts to discrete closure symmetries, appropriate for a tunnelling hierarchy but not proved necessary.
    \item \textbf{Why maximal?} The selection of $2I$ (rather than a smaller group with the same NN angle) via (H4) is motivated by maximal non-resonance but not proved necessary from a variational principle.
    \item \textbf{A dynamical action on the phase bundle.} A field theory on $V \times (S^1)^{|V|}$ with nontrivial kinetic term, under which the six flat coordinates would become genuine Goldstone excitations via Goldstone's theorem. The present paper does not construct such an action.
    \item \textbf{Deriving $\alpha$ from $\Fint$.} Computing the photon coupling coefficient directly from the vertex-exclusion overlap integral, rather than importing it from Paper~XXII.
    \item \textbf{Deriving $\pi\,\lbar_p$ from $\Fint$.} Likewise for the kinematic half-period coefficient~\eqref{eq:kinematic_id}.
    \item \textbf{$O(\alpha^2)$ corrections.} Whether the second-order terms in the perturbative expansion vanish exactly or are merely small.
    \item \textbf{Multi-body extension.} The phase-channel count for nuclei beyond the deuteron.
    \item \textbf{Shell-normalisation hypothesis.} The normalisation $d_2 = 1$ used in Proposition~\ref{prop:phi} to fix absolute scales (and propagated through the $\theta_1 = \pi/5$ angle and the 600-cell selection under H1--H4) is a choice of units, not a derivation. Different shell-normalisation choices rescale $d_1, d_2, \Delta$ simultaneously without changing ratios, but the match to the physical nuclear scale requires a separate calibration.
    \item \textbf{Spectral decomposition of $\Fint$.} The heuristic spectral formula~\eqref{eq:Fint_spectral} is \emph{not} supplied with an exact completeness-relation proof; it is used as a motivation for the two-sector mode structure. An exact spectral identity for the overlap $\eta$ in the Laplacian eigenbasis is an identified open item.
    \item \textbf{Joint operator-space additive decomposition.} The two retained mode sectors (phase-coherence and graph-Laplacian zero mode) live on different operator spaces; we \emph{assume} they combine additively in the free energy without cross-coupling at the order considered. A fibred Hilbert-space / joint-operator construction proving this additivity is an identified open item.
    \item \textbf{Disjoint-support hypothesis for $P_3$ probes.} The six-channel phase-coherence count in Proposition~\ref{prop:pion} uses the assumption that the two three-vertex $P_3$ support graphs are disjoint; vertex-exclusion $i \neq j$ does not by itself force disjoint support. Proving or relaxing this disjointness is an identified open item.
    \item \textbf{Pion-minimum positive-curvature assumption.} Proposition~\ref{thm:factor}'s perturbative algebra assumes $\Fint^{(\pi)}$ has a binding minimum at $r_*^{(\pi)} = 6\,\pi\,\lbar_p$ with positive curvature $K_\pi > 0$. The base $\Fint = -\varepsilon^2\log(1-\eta)$ is repulsive (Proposition~\ref{prop:Fint_props}); the binding pion-sector minimum is therefore a feature of the separate phase-bundle lift, not of $\Fint$ on the base vertex graph. Bridging these --- exhibiting a genuine pion binding minimum from a joint base+fibre construction --- is an identified open item.
    \item \textbf{Classification of finite configurations in $\tfrac{1}{2}\mathbb{Z}[\phig]$.} Remark~\ref{cor:phi_determined} invokes a classification of finite $H_4$-invariant configurations on $S^3$ as an imported standard fact; the analogous classification without $H_4$-invariance is not supplied.
\end{enumerate}

%==================================================================
\section{Reproducibility}\label{sec:reproducibility}
%==================================================================

The derivation chain reduces to:
\begin{enumerate}
    \item $r^2 = r + 1 \Rightarrow r = \phig$, $\cos\theta_1 = \phig/2 \Rightarrow \theta_1 = \pi/5$. (Algebra.)
    \item $2I$ has 120 elements, NN angle $\pi/5$, degree~12. (Finite group theory; verified by explicit quaternion construction in \texttt{build\_600cell.py}.)
    \item $n_\pi = 2 \times 3 = 6$. $\alpha^{-1} = 137 + \pi/87 = 137.0361$. (Arithmetic.)
    \item $r_{\mathrm{str}} = (6 + \alpha)\,\pi \times 0.21031 = 3.96905$~fm. (Calculator.)
    \item $r_d = \sqrt{0.70768 - 0.11762 + 3.93834} = \sqrt{4.52840} = 2.12800$~fm. (Calculator.)
\end{enumerate}

%==================================================================
\section{Conclusion}\label{sec:conclusion}
%==================================================================

The closure functional of the VFD programme is variationally characterised within the stated axiom class rather than merely postulated. Starting from the permeability axiom (self-similar vacuum permeability) --- together with the geometric-mean closure
condition, the shell-normalisation $d_2 = 1$, and a minimal set of structural assumptions (compact
$S^3$, finite-subgroup closure, maximal vertex-transitivity (H1--H4 of Section~\ref{sec:600cell}), and the regularised soft-min axioms A1--A4 of Section~\ref{sec:F}, including A4 Faddeev grouping), the golden ratio is
uniquely selected, the 600-cell is selected from the ADE
classification, the closure functional is uniquely characterised
within its stated class, and the interaction functional follows
from vertex exclusion in the two-basin partition function. Under
the additional structural identifications of phase channels with
pion exchange and of the graph-Laplacian zero mode with photon
exchange (both inherited from prior papers in the
programme~\citep{SmartRadius,SmartXXII}), \emph{together with} the disjoint-support hypothesis of Proposition~\ref{prop:pion} and the pion binding-minimum assumption (H4) of Proposition~\ref{thm:factor}, the resulting
$(6 + \alpha)$ separation factor yields a deuteron charge radius
in agreement with experiment at $0.02\sigma$ against the CODATA 2018 anchor (Pohl~\emph{et~al.}~muonic-deuterium) and $0.83\sigma$ against the CODATA 2022 adjustment, bringing all six charge radii in the programme within $1\sigma$ against both anchors.

The honest classification throughout separates three strata: the
structural spine (derived, $\blacktriangleright$); the physical
identifications that bridge geometry to observables
($\triangleright$); and the open problems ($\circ$), chief among
which are deriving the photon coupling directly from the
inter-basin overlap and constructing a dynamical action on the
phase bundle under which the six flat coordinates become genuine
Goldstone excitations. The framework's permeability axiom --- that the vacuum is a self-similarly permeable field --- is the one physical starting point; everything else (H1--H4, A1--A4, the disjoint-support hypothesis, the pion binding-minimum assumption, and the kinematic identification) is a structural or mathematical assumption whose status is catalogued in Section~\ref{sec:chain}.

%==================================================================
\appendix
\section{Boundary-Resolvent Derivation of $r_p = 4\,\lbar_p$}\label{app:proton}
%==================================================================

For completeness and to make the deuteron calculation self-contained,
we transcribe the derivation of $r_p = 4\,\lbar_p$ from the boundary
resolvent on the proton support graph $P_3$ (full treatment
in~\citep{SmartRadius}).

\subsection*{Proton support graph}

The proton support is $S_p = \{2, 3, 4\}$ (three consecutive 600-cell
shells), inducing the path graph $P_3$ with vertex set $V(P_3) =
\{2, 3, 4\}$ and edges $E(P_3) = \{(2,3), (3,4)\}$. The graph
Laplacian is
\begin{equation}\label{eq:L_P3}
    L(P_3) \;=\; \begin{pmatrix} 1 & -1 & 0 \\ -1 & 2 & -1 \\ 0 & -1 & 1 \end{pmatrix},
    \qquad \mathrm{spec}(L) = \{0, 1, 3\}, \qquad \mathrm{Tr}(L) = 4.
\end{equation}
The eigenvectors are $\psi_0 = \tfrac{1}{\sqrt 3}(1,1,1)^T$ (zero
mode), $\psi_1 = \tfrac{1}{\sqrt 2}(1, 0, -1)^T$ (dipole,
eigenvalue~1), and $\psi_2 = \tfrac{1}{\sqrt 6}(1, -2, 1)^T$
(quadrupole, eigenvalue~3).

\subsection*{Boundary resolvent}

The electromagnetic probe couples at the outer boundary vertex
(shell~4), because interior vertices are exponentially shielded by
closure barriers (see Section~5 of~\citep{SmartRadius}). The
eigenvector weights at vertex 4 are $|\psi_k(4)|^2 = \{1/3, 1/2,
1/6\}$ for eigenvalues $\{0, 1, 3\}$. The boundary resolvent is
\begin{equation}\label{eq:R_bdy}
    R_{\mathrm{bdy}}(z) \;=\; \sum_{k=0}^{2} \frac{|\psi_k(4)|^2}{1 + \lambda_k z}
    \;=\; \frac{1/3}{1} \,+\, \frac{1/2}{1+z} \,+\, \frac{1/6}{1+3z}.
\end{equation}

\subsection*{Charge radius from the form-factor slope}

Identifying $R_{\mathrm{bdy}}(z)$ with the proton electric form factor
$G_E(Q^2)$ at $z = a^2 Q^2$ (with $a^2 = 2\mathrm{Tr}(L)/|V|\,\lbar_p^2$
the heat-kernel cumulant coefficient, Paper~IV~\citep{SmartIV}):
\begin{equation}\label{eq:dR}
    \frac{dR_{\mathrm{bdy}}}{dz}\bigg|_{z=0} \;=\; -\tfrac{1}{2} \,-\, \tfrac{1}{2} \;=\; -1.
\end{equation}
The definition $r_p^2 = -6\,dG_E/dQ^2|_{Q^2 = 0}$ combined with
$dG_E/dQ^2 = a^2\,dR/dz$ gives
\begin{equation}\label{eq:rp_derived}
    r_p^2 \;=\; 6\,a^2 \;=\; 6 \times \frac{2\,\mathrm{Tr}(L)}{|V|}\,\lbar_p^2
    \;=\; 6 \times \frac{8}{3}\,\lbar_p^2 \;=\; 16\,\lbar_p^2,
    \qquad r_p \;=\; 4\,\lbar_p.
\end{equation}

\subsection*{Triple appearance of the factor 4}

The factor $4$ in~\eqref{eq:rp_derived} appears in three related structural characterisations (these are not algebraically independent --- Proposition~\ref{prop:triple} shows items~(1) and~(2) coincide under the assignment rules R1--R5; item~(3) is a heat-kernel reformulation of~(1))~\citep{SmartRadius}:
\begin{enumerate}
    \item $\mathrm{Tr}(L(P_3)) = 2\,|E(P_3)| = 4$ (graph invariant).
    \item $\max(S_p) = 4$ (outermost shell index of the proton support).
    \item $\sqrt{12 \cdot \mathrm{Tr}(L)/|V|} = \sqrt{12 \cdot 4/3} = 4$ (leading heat-kernel cumulant).
\end{enumerate}

\begin{proposition}[Structural identity, from~\citep{SmartRadius}]\label{prop:triple}
For a path graph $P_n$: $\mathrm{Tr}(L) = 2(n-1)$. For a contiguous
shell support $S$ with $n_{\min} = \min S$ and $|S|$ shells:
$\max(S) = n_{\min} + |S| - 1$. These are equal iff
$n_{\min} = |S| - 1$. For the proton support, the assignment rules
R1--R5 of~\citep{SmartRadius} fix $n_{\min} = 2$ and $|S| = 3$,
giving $n_{\min} = 2 = |S| - 1$; hence
$\mathrm{Tr}(L(P_3)) = \max(S_p) = 4$.
\end{proposition}

The boundary-dominance argument (shell-4 vertex is unshielded; interior
vertices are exponentially suppressed by closure barriers of height
$\sim d_1^2/(4\varepsilon^2)$) is verified by the dipole eigenvector
$\psi_1 = (1, 0, -1)/\sqrt{2}$ placing equal and opposite weight on the
two boundary vertices and exactly zero on the interior
(cf.\ Section~4 of~\citep{SmartRadius}).

Numerical evaluation: $\lbar_p = \hbar/(m_p c) = 0.21031$~fm, so
$r_p = 4 \times 0.21031 = 0.84124$~fm, compared to the CODATA value
$0.8409 \pm 0.0004$~fm ($0.04\%$, $0.84\sigma$).

%==================================================================
\bibliographystyle{plainnat}
\bibliography{references}

\end{document}
